\section{Introduction}\label{sec:intro}
A risk certificate must concern all relevant events, including those whose outcomes have not been verified. When harmful outcomes resolve less often, the verified cases need not represent the target population. Even among resolved cases, the time at which a label becomes available may depend on the outcome. A terminal verification indicator discards this second source of information.

The statistical question is therefore more specific than whether timing predicts harm: under a common sensitivity restriction, can the observed resolution path justify a risk decision that the endpoint summary cannot? A useful answer must separate three uncertainties. Unresolved labels create an identification region even with an infinite sample. A sensitivity assumption restricts the completions of those labels. Sampling uncertainty then determines whether a particular risk threshold can be certified during monitoring.

We consider a binary harmful outcome $L$ and a fixed finite resolution path $J^*$, with a distinct state for nonresolution through a terminal horizon. The target is $\theta=\Pbb(L=1)$. We impose $I(L;J^*)\le\rho$, where information is measured in nats and $\rho$ is a declared sensitivity budget. This budget limits average dependence, including dependence through terminal selection. It does not assert missing at random, and it is not estimated from the incompletely labelled stream. Deterministically merging timing bands preserves the same numerical budget by data processing.

Our contribution is an explicit solution for this observation model. Only one terminal harmful mass remains unknown. Its information profile is convex, its minimum is an observable weighted timing information, and its sublevel set gives the sharp risk interval. Refinement increases the profile by a constant independent of the hidden mass. This identity yields exact compatibility and certification frontiers and explains how timing changes decisions. We then project simultaneous categorical confidence sequences through this model, using a numerical outer search whose validity requires conservative enclosures, including when optimization stops before convergence.

The distinction has measurable consequences. A population example changes from an endpoint upper risk of $0.065835$ to $0.055241$ with eight timing bands, crossing a risk budget of $0.06$. A separate, prospectively frozen synthetic evaluation demonstrates finite-sample value under stronger timing contrast: two bands certify 100 of 100 streams by 5,000 matured events, compared with 96 of 100 for the endpoint analysis by 50,000. The matched control quantifies why this benefit depends on timing information rather than on resolution alone.

\section{Related work}\label{sec:related}
\paragraph{Incomplete outcomes and richer observation patterns}
The starting point is the partial-identification approach to missing outcomes: combine the observed distribution with the admissible distributions of missing labels \citep{manski2005}. Sensitivity analysis makes the additional assumptions explicit rather than treating a complete-case estimate as the full-population risk; \citet{carpenter2021} discuss this distinction in practical missing-data analysis. Pattern-mixture models factor distributions by missingness pattern \citep{little1993}. Informative-coarsening sensitivity analysis extends this perspective to partially observed event times \citep{shardell2009}. These are close conceptual precedents for treating the observation process as statistical information. More recently, \citet{dalygrafstein2026} study nonignorable categorical missing outcomes using Bayesian partial identification with additional structural restrictions. Their inferential framework differs from the present mutual-information restriction on the resolution path.

Richer observation types can sharpen bounds. \citet{jiang2018} exploit multiple missingness types under a partial missing-at-random restriction. Callback data can instead support point identification under a stableness-of-resistance assumption linking response mechanisms across contacts \citep{miao2025}. TrajCert uses resolved timing bands and a single unresolved category, with no contact-specific resistance equality or designated missing-at-random type. Its information constraint leaves an interval of terminal completions. The gain from observing the process is established here through an explicit profile difference, not through a general claim that timing data are new.

\paragraph{Sensitivity restrictions and information constraints}
Different sensitivity scales define different admissible laws. \citet{kline2013} bound departures from missing-at-random using a maximal Kolmogorov--Smirnov discrepancy between observed and missing outcome distributions. \citet{christensen2023} optimize counterfactuals over divergence neighborhoods of latent-variable distributions. The $f$-sensitivity model of \citet{jin2026} bounds average selection bias from unmeasured confounding and links treatment-effect bounds to distributionally robust optimization. These results establish substantial precedent for sensitivity analysis based on average discrepancies. Our constraint is the mutual information between the binary outcome and its categorical resolution path; its product-of-marginals reference changes with the unknown outcome marginal. In contrast to the covariate-conditional causal targets of these divergence models, the present observation table leaves one hidden binary split. The contribution is the exact profile, coarsening and decision geometry for that table, combined with conservative sequential projection. Information anchoring in multiple imputation concerns preservation of inferential information about treatment effects \citep{cro2019}, a different meaning of information from Shannon dependence.

\paragraph{Identification and sequential inference}
A confidence region for an identified set and a confidence interval for a particular partially identified parameter have different coverage targets \citep{imbens2004}. \citet{kaido2019} develop calibrated projection confidence intervals for partially identified parameters under moment restrictions. Sequential inference for partially identified causal sensitivity parameters is also considered by \citet{dalal2024} under asymptotic conditions. Our construction instead projects a finite-sample categorical confidence sequence through the identified set induced by the resolution-path information budget. \citet{lindon2022} provide joint anytime-valid inference for multinomial data, offering an alternative to the coordinatewise categorical construction used here. Our construction covers the complete partition-specific identified set from above, uniformly in monitoring time. Confidence sequences and mixture martingales are established tools \citep{howard2021,ramdas2023}; betting methods provide further constructions for bounded means \citep{waudby2024}, and always-valid inference supports continuous monitoring of experiments \citep{johari2022}. Delayed-outcome inference also requires care about observation order: the recent design-based work of \citet{lindon2026} targets cumulative counterfactual rewards under staggered entry and outcome delay. We instead target a stable population risk with potentially unresolved labels, admitting each event only at a common terminal age. The sequential contribution is the sensitivity projection and its numerical implementation, conditional on conservative computation and using an established confidence-sequence input. Neither data processing nor martingale inversion is claimed as a new general principle.
